% ---------------------------------------------------------------------
\section{L’impression à la demande, c’est quoi ?}
% ---------------------------------------------------------------------

\subsection{Le principe : imprimer après la vente}

Dans le modèle classique, l’éditeur fait imprimer un \emph{tirage}
(500, 1\,000, 3\,000 exemplaires…), le stocke chez un distributeur,
puis espère le vendre. Les invendus reviennent (les \emph{retours}) et
finissent parfois au pilon.

L’impression à la demande renverse la logique : \textbf{le livre
n’existe que sous forme de fichiers} (un PDF intérieur, un PDF de
couverture) tant que personne ne l’a acheté. Quand un lecteur ou un
libraire passe commande, l’exemplaire est imprimé — souvent à l’unité —
puis expédié, en général sous quelques jours.

\begin{center}
\begin{tikzpicture}[
  etape/.style={draw=plOrange,line width=1.2pt,fill=plCreme,
    minimum width=2.9cm,minimum height=1.2cm,align=center,
    font=\small\bfseries,text width=2.7cm},
  fleche/.style={-latex,line width=1.6pt,plNoir}]
  \node[etape] (a) {Commande\\du lecteur};
  \node[etape,right=0.7cm of a] (b) {Plateforme\\ou distributeur};
  \node[etape,right=0.7cm of b] (c) {Impression\\numérique};
  \node[etape,right=0.7cm of c] (d) {Expédition\\(lecteur ou libraire)};
  \draw[fleche] (a) -- (b);
  \draw[fleche] (b) -- (c);
  \draw[fleche] (c) -- (d);
\end{tikzpicture}
\end{center}

\subsection{Trois façons d’imprimer un livre}

\begin{center}
\small
\begin{tabularx}{\linewidth}{P{3.1cm}LLL}
\entete & \textbf{Offset} & \textbf{Numérique petit tirage} & \textbf{Impression à la demande} \\
Quantité typique & à partir de 500 – 1\,000 ex. & 30 à 500 ex. & 1 ex. à la commande \\
Coût unitaire & le plus bas à fort volume & moyen & le plus élevé \\
Mise de départ & élevée (plaques, calage) & modérée & quasi nulle \\
Stock & oui, à financer et stocker & petit stock & aucun \\
Délai & 2 à 4 semaines & 1 à 2 semaines & quelques jours par commande \\
Qualité & référence, tous papiers & très bonne & bonne à très bonne, choix de papiers limité \\
Idéal pour & best-sellers, beaux livres & lancements, salons & fonds, tests, niches, autoédition \\
\bottomrule
\end{tabularx}
\end{center}

\subsection{Pourquoi c’est intéressant pour un petit éditeur ou un auteur}

\begin{itemize}
  \item \textbf{Aucun stock à financer} : la trésorerie n’est pas
        immobilisée, pas de frais d’entrepôt, pas de pilon.
  \item \textbf{Un livre jamais épuisé} : comme il n’y a pas de stock
        qui s’écoule, le livre reste en vente aussi longtemps que vous
        le souhaitez. Même un titre ancien qui ne se vend qu’à quelques
        exemplaires par an continue de rapporter, sans rien coûter à
        stocker.
  \item \textbf{Un risque limité} : on peut tester un titre, un genre,
        une couverture, et corriger le fichier à tout moment.
  \item \textbf{Un accès direct aux grands canaux} : Amazon, Fnac,
        Cultura, Decitre… et, selon la plateforme, les librairies
        indépendantes via Dilicom.
\end{itemize}

\subsection{Les limites à connaître}

\begin{itemize}
  \item \textbf{Une marge unitaire plus faible} qu’en offset : chaque
        exemplaire coûte plus cher à fabriquer.
  \item \textbf{Une présence en librairie difficile} : les livres en IAD
        sont commandables, mais rarement en rayon. Sans diffuseur
        (représentants qui démarchent les libraires), le livre ne se
        vend pas « tout seul ».
  \item \textbf{Des finitions limitées} : formats standards, quelques
        papiers, peu ou pas de dorure, vernis sélectif ou rabats.
  \item \textbf{Une qualité qui peut varier} : selon l’imprimerie qui
        fabrique l’exemplaire, les couleurs, le centrage du texte ou la
        coupe des pages peuvent légèrement différer. Avant de lancer la
        vente, commandez toujours un exemplaire test (une « épreuve »)
        pour vérifier le résultat.
\end{itemize}

\begin{astuce}[En résumé]
L’IAD n’est pas « l’édition au rabais » : c’est un \textbf{outil de
fabrication et de logistique}. Beaucoup d’éditeurs combinent les deux :
un petit tirage pour les salons et les librairies qui suivent le titre,
l’IAD pour assurer la disponibilité permanente partout ailleurs.
\end{astuce}
